\documentclass{article}
\usepackage{amsmath,amssymb,amsthm}
\usepackage{algorithm,algorithmic}
\usepackage{hyperref}
\newtheorem{theorem}{Theorem}
\newtheorem{lemma}{Lemma}
\newtheorem{assumption}{Assumption}
\newtheorem{remark}{Remark}
\newcommand{\prox}{\operatorname{prox}}
\newcommand{\R}{\mathbb{R}}
\newcommand{\E}{\mathbb{E}}

\begin{document}

\section*{Proximal Gradient Method (Forward–Backward Splitting)}

\section*{Mathematical Formulation}
Let $f:\R^{d}\rightarrow\R$ be differentiable and $g:\R^{d}\rightarrow\R\cup\{+\infty\}$ be a proper, closed, convex (possibly non‑smooth) function.  
Define the composite objective
\[
\Phi(x)\;:=\;f(x)+g(x),\qquad x\in\R^{d}.
\]

Assume that $\nabla f$ is $L$‑Lipschitz continuous, i.e.
\[
\|\nabla f(x)-\nabla f(y)\|\le L\|x-y\|,\qquad\forall x,y\in\R^{d}.
\]

Given a stepsize $\eta\in(0,1/L]$, the proximal gradient iteration is
\begin{equation}
\label{eq:pg}
x_{k+1}\;=\;\prox_{\eta g}\!\bigl(x_{k}-\eta \nabla f(x_{k})\bigr),\qquad k=0,1,2,\dots
\end{equation}
where the proximal operator of $g$ with parameter $\eta$ is
\[
\prox_{\eta g}(u)\;:=\;\arg\min_{z\in\R^{d}}\Bigl\{ g(z)+\frac{1}{2\eta}\|z-u\|^{2}\Bigr\}.
\]

\section*{Pseudocode}
\begin{algorithm}[H]
\caption{Proximal Gradient Method}
\begin{algorithmic}[1]
\STATE \textbf{Input:} stepsize $\eta\in(0,1/L]$, initial point $x_{0}\in\R^{d}$, max iter $K$.
\FOR{$k=0$ \TO $K-1$}
    \STATE Compute gradient $g_{k}\gets \nabla f(x_{k})$.
    \STATE Gradient step $y_{k}\gets x_{k}-\eta g_{k}$.
    \STATE Proximal step $x_{k+1}\gets \prox_{\eta g}(y_{k})$.
\ENDFOR
\STATE \textbf{Output:} $x_{K}$ (or any iterate $x_{k}$).
\end{algorithmic}
\end{algorithm}

\section*{Dry Run Example}
Consider the scalar problem
\[
f(w)=(w-3)^{2},\qquad g(w)=0,\qquad \eta=0.1,\qquad w_{0}=0.
\]
Since $g\equiv0$, $\prox_{\eta g}(u)=u$.

\begin{center}
\begin{tabular}{c|c|c|c}
$k$ & $w_{k}$ & $\nabla f(w_{k})$ & $w_{k+1}=w_{k}-\eta\nabla f(w_{k})$ \\ \hline
0 & $0$ & $2(w_{0}-3)=-6$ & $0-0.1(-6)=0.6$ \\
1 & $0.6$ & $2(0.6-3)=-4.8$ & $0.6-0.1(-4.8)=1.08$ \\
2 & $1.08$ & $2(1.08-3)=-3.84$ & $1.08-0.1(-3.84)=1.464$ \\
\end{tabular}
\end{center}
Objective values:
\[
\Phi(w_{0})=9,\quad \Phi(w_{1})=(0.6-3)^{2}=5.76,\quad \Phi(w_{2})=(1.08-3)^{2}=3.6864,\quad \Phi(w_{3})=(1.464-3)^{2}=2.366.
\]
The iterates approach the minimiser $w^{\star}=3$.

\newpage
\section*{Assumptions}
\begin{assumption}[Convexity and Smoothness]\label{as:convex}
\begin{enumerate}
\item $f$ is convex and differentiable on $\R^{d}$.
\item $\nabla f$ is $L$‑Lipschitz continuous for some $L>0$.
\item $g$ is proper, closed and convex (possibly non‑smooth).
\end{enumerate}
\end{assumption}
\begin{assumption}[Stepsize]\label{as:stepsize}
The stepsize satisfies $0<\eta\le 1/L$.
\end{assumption}

\section*{Main Convergence Theorem}
\begin{theorem}[Ergodic $O(1/k)$ rate]\label{thm:rate}
Let $\{x_{k}\}_{k\ge0}$ be generated by \eqref{eq:pg} under Assumptions~\ref{as:convex}–\ref{as:stepsize}.  
Define the averaged iterate
\[
\bar{x}_{K}:=\frac{1}{K}\sum_{k=0}^{K-1}x_{k}.
\]
Then for any optimal solution $x^{\star}\in\arg\min\Phi$,
\begin{equation}
\label{eq:rate}
\Phi(\bar{x}_{K})-\Phi(x^{\star})\;\le\;\frac{\|x_{0}-x^{\star}\|^{2}}{2\eta K}.
\end{equation}
Consequently $\Phi(\bar{x}_{K})\to\Phi(x^{\star})$ at the rate $O(1/K)$.
\end{theorem}

\section*{Theoretical Properties}
\begin{itemize}
\item \textbf{Stability.} Because the proximal operator is non‑expansive, the iteration \eqref{eq:pg} is a contraction in the metric induced by the Euclidean norm, guaranteeing boundedness of the whole sequence.
\item \textbf{Hyperparameter Sensitivity.} The stepsize must respect $\eta\le1/L$; larger $\eta$ may violate the descent property derived in the proof.
\item \textbf{Comparison.} When $g\equiv0$, \eqref{eq:pg} reduces to gradient descent and the bound \eqref{eq:rate} coincides with the classic $O(1/k)$ result. When $g$ is the $\ell_{1}$ norm, the method becomes the Iterative Soft‑Thresholding Algorithm (ISTA) with the same rate.
\end{itemize}

\section*{Discussion}
The proof hinges on two elementary facts:
\begin{enumerate}
\item \emph{Descent Lemma} for $L$‑smooth $f$.
\item \emph{Non‑expansiveness} of the proximal operator:
\[
\|\prox_{\eta g}(u)-\prox_{\eta g}(v)\|\le\|u-v\|,\qquad\forall u,v\in\R^{d}.
\]
These combine to produce a one‑step inequality that telescopes when summed over iterations, yielding the $O(1/K)$ bound.

\newpage
\section*{Preliminaries (Definitions and Lemmas)}
\begin{lemma}[Descent Lemma]\label{lem:descent}
If $\nabla f$ is $L$‑Lipschitz, then for any $x,y\in\R^{d}$,
\[
f(y)\le f(x)+\langle\nabla f(x),y-x\rangle+\frac{L}{2}\|y-x\|^{2}.
\]
\end{lemma}
\begin{proof}
Standard proof by integrating the Lipschitz condition. \qedhere
\end{proof}

\begin{lemma}[Firm non‑expansiveness of the proximal operator]\label{lem:firm}
For any $u,v\in\R^{d}$,
\[
\|\prox_{\eta g}(u)-\prox_{\eta g}(v)\|^{2}
\le\langle u-v,\prox_{\eta g}(u)-\prox_{\eta g}(v)\rangle .
\]
\end{lemma}
\begin{proof}
Follows from the optimality condition of the proximal subproblem; see e.g. \cite{parikh2014proximal}. \qedhere
\end{proof}

\section*{Main Proof (Step‑by‑Step Derivation)}
We prove Theorem~\ref{thm:rate}. Throughout the proof let $x^{\star}\in\arg\min\Phi$ be arbitrary but fixed.

\paragraph{Step 1 (One‑step inequality).}
From the definition \eqref{eq:pg} set
\[
y_{k}=x_{k}-\eta\nabla f(x_{k}),\qquad x_{k+1}=\prox_{\eta g}(y_{k}).
\]
Apply Lemma~\ref{lem:firm} with $u=y_{k}$ and $v=x^{\star}$ (note that $x^{\star}=\prox_{\eta g}(x^{\star}-\eta\nabla f(x^{\star})$ because $0\in\partial g(x^{\star})+\nabla f(x^{\star})$):
\begin{align}
\|x_{k+1}-x^{\star}\|^{2}
&\le\langle y_{k}-x^{\star},\,x_{k+1}-x^{\star}\rangle \nonumber\\
&=\langle x_{k}-\eta\nabla f(x_{k})-x^{\star},\,x_{k+1}-x^{\star}\rangle . \tag{Step 1}
\end{align}

\paragraph{Step 2 (Expand the inner product).}
Add and subtract $x_{k}$ inside the second factor:
\[
\langle x_{k}-x^{\star},\,x_{k+1}-x^{\star}\rangle
   -\eta\langle \nabla f(x_{k}),\,x_{k+1}-x^{\star}\rangle .\tag{Step 2}
\]

\paragraph{Step 3 (Use Cauchy–Schwarz and Young).}
For any vectors $a,b$ and $\alpha>0$, $2\langle a,b\rangle\le\alpha\|a\|^{2}+ \alpha^{-1}\|b\|^{2}$.  
Apply with $a=x_{k+1}-x_{k}$ and $b=x_{k}-x^{\star}$, $\alpha=1$:
\[
\langle x_{k}-x^{\star},\,x_{k+1}-x^{\star}\rangle
= \tfrac12\bigl(\|x_{k}-x^{\star}\|^{2}+\|x_{k+1}-x^{\star}\|^{2}
               -\|x_{k+1}-x_{k}\|^{2}\bigr). \tag{Step 3}
\]

\paragraph{Step 4 (Bound the gradient term).}
Since $x^{\star}$ minimizes $\Phi$, the optimality condition gives
\[
0\in\nabla f(x^{\star})+\partial g(x^{\star})\quad\Longrightarrow\quad
-\nabla f(x^{\star})\in\partial g(x^{\star}).
\]
Using convexity of $g$,
\[
g(x_{k+1})\ge g(x^{\star})+\langle s^{\star},x_{k+1}-x^{\star}\rangle ,
\qquad s^{\star}\in\partial g(x^{\star})=-\nabla f(x^{\star}).
\]
Hence
\[
\langle \nabla f(x^{\star}),x_{k+1}-x^{\star}\rangle\ge g(x^{\star})-g(x_{k+1}). \tag{Step 4}
\]

\paragraph{Step 5 (Combine Steps 1–4).}
Insert the decomposition \eqref{Step 3} into \eqref{Step 2} and rearrange:
\[
\|x_{k+1}-x^{\star}\|^{2}
\le \|x_{k}-x^{\star}\|^{2}
      -\|x_{k+1}-x_{k}\|^{2}
      -2\eta\langle \nabla f(x_{k})-\nabla f(x^{\star}),\,x_{k+1}-x^{\star}\rangle . \tag{Step 5}
\]

\paragraph{Step 6 (Apply the Descent Lemma).}
From Lemma~\ref{lem:descent} with $x=x_{k}$, $y=x_{k+1}$,
\[
f(x_{k+1})\le f(x_{k})+\langle\nabla f(x_{k}),x_{k+1}-x_{k}\rangle
            +\frac{L}{2}\|x_{k+1}-x_{k}\|^{2}. \tag{Step 6}
\]
Rearrange to obtain
\[
\langle\nabla f(x_{k}),x_{k+1}-x_{k}\rangle
\ge f(x_{k+1})-f(x_{k})-\frac{L}{2}\|x_{k+1}-x_{k}\|^{2}. \tag{Step 6$'$}
\]

\paragraph{Step 7 (Eliminate $\|x_{k+1}-x_{k}\|^{2}$).}
Because $\eta\le 1/L$, we have $L\eta\le1$ and therefore
\[
\|x_{k+1}-x_{k}\|^{2}
= \|\eta\nabla f(x_{k})+\bigl(x_{k+1}-y_{k}\bigr)\|^{2}
\ge \eta^{2}\|\nabla f(x_{k})\|^{2}
    -2\eta\langle\nabla f(x_{k}),x_{k+1}-y_{k}\rangle .
\]
Using the optimality condition of the proximal step,
\[
\frac{1}{\eta}(y_{k}-x_{k+1})\in\partial g(x_{k+1}) ,
\]
hence
\[
\langle \nabla f(x_{k}),x_{k+1}-y_{k}\rangle
= -\eta\langle \nabla f(x_{k}),\,\partial g(x_{k+1})\rangle
\le \eta\bigl(g(x_{k})-g(x_{k+1})\bigr)    \tag{Step 7}
\]
by convexity of $g$.

\paragraph{Step 8 (Putting everything together).}
Insert \eqref{Step 6$'$} and \eqref{Step 7} into \eqref{Step 5} and use $\eta\le1/L$ to cancel the $\|x_{k+1}-x_{k}\|^{2}$ terms. After straightforward algebra we obtain
\[
\Phi(x_{k+1})-\Phi(x^{\star})
\le \frac{1}{2\eta}\Bigl(\|x_{k}-x^{\star}\|^{2}
                         -\|x_{k+1}-x^{\star}\|^{2}\Bigr). \tag{Step 8}
\]

\paragraph{Step 9 (Telescoping sum).}
Sum the inequality in \eqref{Step 8} for $k=0,\dots,K-1$:
\[
\sum_{k=0}^{K-1}\bigl[\Phi(x_{k+1})-\Phi(x^{\star})\bigr]
\le \frac{1}{2\eta}\bigl(\|x_{0}-x^{\star}\|^{2}
                       -\|x_{K}-x^{\star}\|^{2}\bigr)
\le \frac{\|x_{0}-x^{\star}\|^{2}}{2\eta}. \tag{Step 9}
\]

\paragraph{Step 10 (Convexity of $\Phi$ and averaging).}
Since $\Phi$ is convex,
\[
\Phi(\bar{x}_{K})-\Phi(x^{\star})
\le\frac{1}{K}\sum_{k=1}^{K}\bigl[\Phi(x_{k})-\Phi(x^{\star})\bigr].
\]
Combine with \eqref{Step 9} (shift index by one) to obtain exactly the bound \eqref{eq:rate}.

\paragraph{Conclusion.}
Thus Theorem~\ref{thm:rate} holds. \qed

\section*{Rate Derivation (With Constants)}
The constant in \eqref{eq:rate} is $\frac{1}{2\eta}$.
Choosing the largest admissible stepsize $\eta=1/L$ yields the tightest bound
\[
\Phi(\bar{x}_{K})-\Phi(x^{\star})\le \frac{L}{2}\,\frac{\|x_{0}-x^{\star}\|^{2}}{K}.
\]

\section*{Special Cases}
\begin{itemize}
\item \textbf{Pure gradient descent} ($g\equiv0$): the proximal step disappears, and the proof reduces to the classical $O(1/K)$ rate for smooth convex minimisation.
\item \textbf{Lasso (elastic net)}: $g(w)=\lambda\|w\|_{1}$. The proximal operator becomes soft‑thresholding; the same $O(1/K)$ guarantee applies to ISTA.
\item \textbf{Strongly convex $f+g$}: if $\Phi$ is $\mu$‑strongly convex, a similar derivation (adding the strong convexity term) yields a linear rate $O\bigl((1-\eta\mu)^{K}\bigr)$.
\end{itemize}

\begin{thebibliography}{9}
\bibitem{parikh2014proximal}
N. Parikh and S. Boyd, \emph{Proximal Algorithms}, Foundations and Trends in Optimization, 2014.
\end{thebibliography}

\end{document}